\documentclass[10pt,letterpaper]{article}
\usepackage{cornell-notes}

\title{Clause 06.08: Inclusion Of View Components}
\author{}
\date{}
\setDocTitle{Clause 06.08: Inclusion Of View Components}
\setDocSubtitle{ISO/IEC/IEEE 42010:2022}
\setDocOwner{ISO 42010 Cornell Notes}
\setDocDate{\today}
\setCornellCollection{ISO/IEC/IEEE 42010:2022 Cornell Notes}
\setCornellUnitType{Clause}
\setCornellUnitNumber{06.08}
\setCornellUnitTitle{Inclusion Of View Components}
\hypersetup{pdftitle={Clause 06.08: Inclusion Of View Components},pdfsubject={Cornell notes},pdfauthor={}}

\begin{document}
\maketitle
\begin{CornellOverviewBox}{Overview}
{\Large\bfseries\textcolor{CNavy}{Section 6.8: Inclusion of view components}}\\[3pt]
{\small\textbf{Classification:} Normative}\\
{\small\textbf{Learning objective:} Identify the mandatory, recommended, and permitted provisions governing inclusion of view components.}
\end{CornellOverviewBox}
\vspace{3pt}

\begin{CornellNotesTable}
\CornellNoteRow{Purpose}{Establishes how view components compose views and how model kinds or legends govern their conventions.}
\CornellNoteRow{Core details}{\begin{itemize}[leftmargin=*,nosep,topsep=1pt]
\item Within a view, one or more view components can be used to selectively present some or all of the informational content required by the architecture viewpoint to highlight points of interest.
\item Correspondences can express the relationships among components shared across views.
\end{itemize}}
\CornellNoteRow{Requirement}{\begin{itemize}[leftmargin=*,nosep,topsep=1pt]
\item An architecture view must be composed of one or more view components in accordance with its governing architecture viewpoint.
\item Each view component must include version identification as specified by the organization and/or project.
\item Each view component must identify its governing model kind, if any, and adhere to the conventions of its governing architecture viewpoint (see 6.6).
\item When a view component does not have a governing model kind, i.e. for an information part not described with a model, a view component legend must be included to specify the conventions used in that view component.
\end{itemize}}
\CornellNoteRow{Permission}{\begin{itemize}[leftmargin=*,nosep,topsep=1pt]
\item A view component may be a part of more than one architecture view.
\end{itemize}}
\CornellNoteRow{Example / application}{A view component that is a record of expert opinion, rather than a model that one can analyse using calculations, simulation, or any other suitable analysis method.}
\CornellNoteRow{Interpretive note}{\begin{itemize}[leftmargin=*,nosep,topsep=1pt]
\item Sharing view components between architecture views permits an AD to capture distinct but related concerns without redundancy or repetition of the same information in multiple views and reduces possibilities for inconsistency.
\item No particular method is prescribed for the level of formality of view components to be used in an AD. While model-based view components that have a formal specification of semantics and syntax can be less ambiguous, non-model-based view components can also be used effectively.
\end{itemize}}
\CornellNoteRow{Related sections}{6.6}
\CornellNoteRow{Conformance check}{\begin{itemize}[leftmargin=*,nosep,topsep=1pt]
\item An architecture view must be composed of one or more view components in accordance with its governing architecture viewpoint
\item Each view component must include version identification as specified by the organization and/or project
\end{itemize}}
\CornellNoteRow{Active recall}{\begin{itemize}[leftmargin=*,nosep,topsep=1pt]
\item Which provisions in Inclusion of view components are mandatory for conformance?
\item Which provisions are recommendations or permissions rather than requirements?
\end{itemize}}
\end{CornellNotesTable}


\begin{CornellSummaryBox}{Summary}
Establishes how view components compose views and how model kinds or legends govern their conventions. For conformance, its mandatory provisions must be satisfied, including: An architecture view must be composed of one or more view components in accordance with its governing architecture viewpoint; Each view component must include version identification as specified by the organization and/or project. The section also explicitly permits: A view component may be a part of more than one architecture view.
\end{CornellSummaryBox}

\end{document}
